\documentclass{beamer}
\mode<presentation>
\usepackage{amsmath,amssymb,mathtools}
\usepackage{textcomp}
\usepackage{gensymb}
\usepackage{adjustbox}
\usepackage{subcaption}
\usepackage{enumitem}
\usepackage{multicol}
\usepackage{listings}
\usepackage{url}
\usepackage{graphicx} % <-- needed for images
\def\UrlBreaks{\do\/\do-}

\usetheme{Boadilla}
\usecolortheme{lily}
\setbeamertemplate{footline}{
  \leavevmode%
  \hbox{%
  \begin{beamercolorbox}[wd=\paperwidth,ht=2ex,dp=1ex,right]{author in head/foot}%
    \insertframenumber{} / \inserttotalframenumber\hspace*{2ex}
  \end{beamercolorbox}}%
  \vskip0pt%
}
\setbeamertemplate{navigation symbols}{}

\lstset{
  frame=single,
  breaklines=true,
  columns=fullflexible,
  basicstyle=\ttfamily\tiny   % tiny font so code fits
}

\numberwithin{equation}{section}

% ---- your macros ----
\providecommand{\nCr}[2]{\,^{#1}C_{#2}}
\providecommand{\nPr}[2]{\,^{#1}P_{#2}}
\providecommand{\mbf}{\mathbf}
\providecommand{\pr}[1]{\ensuremath{\Pr\left(#1\right)}}
\providecommand{\qfunc}[1]{\ensuremath{Q\left(#1\right)}}
\providecommand{\sbrak}[1]{\ensuremath{{}\left[#1\right]}}
\providecommand{\lsbrak}[1]{\ensuremath{{}\left[#1\right.}}
\providecommand{\rsbrak}[1]{\ensuremath{\left.#1\right]}}
\providecommand{\brak}[1]{\ensuremath{\left(#1\right)}}
\providecommand{\lbrak}[1]{\ensuremath{\left(#1\right.}}
\providecommand{\rbrak}[1]{\ensuremath{\left.#1\right)}}
\providecommand{\cbrak}[1]{\ensuremath{\left\{#1\right\}}}
\providecommand{\lcbrak}[1]{\ensuremath{\left\{#1\right.}}
\providecommand{\rcbrak}[1]{\ensuremath{\left.#1\right\}}}
\theoremstyle{remark}
\newtheorem{rem}{Remark}
\newcommand{\sgn}{\mathop{\mathrm{sgn}}}
\providecommand{\abs}[1]{\left\vert#1\right\vert}
\providecommand{\res}[1]{\Res\displaylimits_{#1}}
\providecommand{\norm}[1]{\lVert#1\rVert}
\providecommand{\mtx}[1]{\mathbf{#1}}
\providecommand{\mean}[1]{E\left[ #1 \right]}
\providecommand{\fourier}{\overset{\mathcal{F}}{ \rightleftharpoons}}
\providecommand{\system}{\overset{\mathcal{H}}{ \longleftrightarrow}}
\providecommand{\dec}[2]{\ensuremath{\overset{#1}{\underset{#2}{\gtrless}}}}
\newcommand{\myvec}[1]{\ensuremath{\begin{pmatrix}#1\end{pmatrix}}}
\let\vec\mathbf

\title{Matgeo Presentation - Problem 12.602}
\author{ee25btech11063 - Vejith}

\begin{document}


\frame{\titlepage}
\begin{frame}{Question}
If the following system has a non-trivial solution:
\begin{align*}
px + qy + rz &= 0, \\
qx + ry + pz &= 0,\\
rx + py + qz &= 0,
\end{align*}
then which one of the following options is \textbf{TRUE}? \hfill (CS 2016)
\vspace{0.5cm}

\begin{enumerate}
    \item[(a)] $p + q + r = 0$ \text{ or } $p = q = r$
    \item[(b)] $p + q + r = 0$ \text{ or } $p = q = -r$
    \item[(c)] $p + q + r = 0$ \text{ or } $p = q = r = 0$
    \item[(d)] $p + q + r = 0$ \text{ or } $p = -q = r$
\end{enumerate}
\end{frame}

\begin{frame}{Solution}
    Given equations are
\begin{align}
 \brak{p\hspace{0.5cm}q\hspace{0.5cm}r}\myvec{x\\y\\z}=0\\
   \brak{q\hspace{0.5cm}r\hspace{0.5cm}p}\myvec{x\\y\\z}=0\\
   \brak{r\hspace{0.5cm}p\hspace{0.5cm}q}\myvec{x\\y\\z}=0\\
   \implies \begin{pmatrix}
        p & q & r\\
        q & r & p\\
        r & p & q
    \end{pmatrix} \myvec{x\\y\\z} =\myvec{0\\0\\0}
\end{align}
\end{frame}
\begin{frame}{Solution}
let 
\begin{align}
    \Vec{A}=\begin{pmatrix}
        p & q & r\\
        q & r & p\\
        r & p & q
    \end{pmatrix} \text{ and } \Vec{X}=\myvec{x\\y\\z}\\
    \implies \Vec{A}\Vec{X}=\Vec{0}
\end{align}
The system of equations $\Vec{AX = 0}$ has non trivial solution if rank of $\Vec{A}<3$
\begin{align}
      \begin{pmatrix}
        p & q & r\\
        q & r & p\\
        r & p & q
    \end{pmatrix} &\xleftrightarrow{R_2 \rightarrow pR_2- qR_1} 
    \begin{pmatrix}
        p & q & r \\
        0 & pr-q^2 & p^2-qr \\
        r & p & q 
\end{pmatrix}\\
&\xleftrightarrow{R_3 \rightarrow pR_3- rR_1} \begin{pmatrix}
        p & q & r \\
        0 & pr-q^2 & p^2-qr \\
        0 & p^2-rq & pq-r^2 
\end{pmatrix}\\
&\xleftrightarrow{R_3 \rightarrow (pr-q^2)R_3 - (p^2-rq)R_2} \begin{pmatrix}
        p & q & r \\
        0 & pr-q^2 & p^2-qr \\
        0 & 0 & D 
\end{pmatrix}
\end{align}
\end{frame}

\begin{frame}{Solution}
\begin{align}
\implies D=(pr-q^2)(pq-r^2)-(p^2-rq)(p^2-qr)=0\\
\implies (pr-q^2)(pq-r^2)=p^2qr-pr^3-pq^3+q^2r^2\\
\implies (p^2-rq)(p^2-qr)=p^4-2p^2qr+r^2q^2
\end{align}
From (0.11) and (0.12)
\begin{align}
    D=\brak{p^2qr-pr^3-pq^3+q^2r^2}-\brak{p^4-2p^2qr+r^2q^2}\\
    =-p^4+3p^2qr-pq^3-pr^3\\
    =-p\brak{p^3+q^3+r^3-3pqr}\\
    =-p\brak{p+q+r}\brak{p^2+q^2+r^2-pq-qr-pr}\\
 \implies D  =-\frac{1}{2}p\brak{p+q+r}\brak{\brak{p-q}^2 + \brak{q-r}^2 + \brak{p-r}^2}
\end{align}
$D=0$ means 
\begin{align}
    p+q+r=0 \hspace{0.5cm} \text{ or }\hspace{0.5cm} \brak{p-q}^2 + \brak{q-r}^2 + \brak{p-r}^2=0\\
    \implies  p+q+r=0 \hspace{1cm} \text{ or }\hspace{1cm} p=q=r
    \end{align}
\end{frame}

\begin{frame}{Conclusion}
 $\implies$ Given system of equations has non trivial solution if
 \begin{align}
     p+q+r=0 \hspace{1cm} \text{ or }\hspace{1cm} p=q=r
 \end{align}
\end{frame}


 % --------- CODE APPENDIX ---------
\section*{Appendix: Code}

% C program
\begin{frame}[fragile]{C Code: program.c}
\begin{lstlisting}[language=C]
#include <stdio.h>
#include <math.h>

int main() {
    FILE *fp;

    // Open output file
    fp = fopen("answer.dat", "w");
    if (fp == NULL) {
        printf("Error opening file!\n");
        return 1;
    }

    // ---------------- Example 1: p + q + r = 0 ----------------
    double p1 = 1, q1 = 2, r1 = -3; // satisfies p + q + r = 0
    fprintf(fp, "Example 1:\n");
    fprintf(fp, "Given values: p = %.2lf, q = %.2lf, r = %.2lf\n", p1, q1, r1);

    if (fabs(p1 + q1 + r1) < 1e-9 || (fabs(p1 - q1) < 1e-9 && fabs(q1 - r1) < 1e-9)) {
        fprintf(fp, "→ The system has a NON-TRIVIAL solution.\n");
        if (fabs(p1 + q1 + r1) < 1e-9)
            fprintf(fp, "  Condition satisfied: p + q + r = 0\n\n");
        else
            fprintf(fp, "  Condition satisfied: p = q = r\n\n");
    } else {
        fprintf(fp, "→ The system does NOT have a non-trivial solution.\n\n");
    }
\end{lstlisting}
\end{frame}

\begin{frame}[fragile]{C Code: program.c}
\begin{lstlisting}[language=C]
    // ---------------- Example 2: p = q = r ----------------
    double p2 = 5, q2 = 5, r2 = 5; // satisfies p = q = r
    fprintf(fp, "Example 2:\n");
    fprintf(fp, "Given values: p = %.2lf, q = %.2lf, r = %.2lf\n", p2, q2, r2);

    if (fabs(p2 + q2 + r2) < 1e-9 || (fabs(p2 - q2) < 1e-9 && fabs(q2 - r2) < 1e-9)) {
        fprintf(fp, "→ The system has a NON-TRIVIAL solution.\n");
        if (fabs(p2 + q2 + r2) < 1e-9)
            fprintf(fp, "  Condition satisfied: p + q + r = 0\n");
        else
            fprintf(fp, "  Condition satisfied: p = q = r\n");
    } else {
        fprintf(fp, "→ The system does NOT have a non-trivial solution.\n");
    }

    fclose(fp);
    printf("Results for both examples written to answer.dat\n");
    return 0;
}

\end{lstlisting}
\end{frame}

% Python plotting
\begin{frame}[fragile]{Python: plot.py}
\begin{lstlisting}[language=Python]
import numpy as np

# Example 1: satisfies p + q + r = 0
p, q, r = 1, 2, -3

# Example 2 (uncomment to test p=q=r case)
# p, q, r = 5, 5, 5

# Construct coefficient matrix A
A = np.array([
    [p, q, r],
    [q, r, p],
    [r, p, q]
], dtype=float)

# Compute determinant
detA = np.linalg.det(A)

print(f"Given: p={p}, q={q}, r={r}")
print("Matrix A =\n", A)
print(f"Determinant of A = {detA:.6f}")

# Check for non-trivial solution
if abs(detA) < 1e-9:
    print("\n The system has a NON-TRIVIAL solution.")
    if abs(p + q + r) < 1e-9:
        print("Condition satisfied: p + q + r = 0")
    elif abs(p - q) < 1e-9 and abs(q - r) < 1e-9:
        print("Condition satisfied: p = q = r")
    else:
        print("Condition satisfied: determinant = 0 (degenerate case)")
else:
    print("\n The system has only the TRIVIAL solution (x = y = z = 0).")

\end{lstlisting}
\end{frame}

\end{document}
